\chapter{... que ressuscitou ao terceiro dia; subiu aos céus, está sentado à direita de Deus Pai todo-poderoso}
\chaptermark{... que ressuscitou ao terceiro dia}

Este capítulo aprofunda a continuação do segundo artigo do Símbolo Apostólico, professando a fé na Ressurreição, Ascensão e reinado glorioso de Jesus Cristo. Desdobramos os parágrafos correspondentes do Catecismo da Igreja Católica (CIC), revelando o evento central da fé, a entrada definitiva no céu e o senhorio escatológico de Cristo. Para adultos em processo de iniciação cristã pelo RICA, esses mistérios pascais iluminam a esperança na Vigília Pascal, unindo morte à vida eterna.

\section{A Ressurreição de Jesus Cristo (\S 638--658)}

A Ressurreição é o coroamento da fé cristã, verdade central acreditada pela primeira comunidade, transmitida pela Tradição e confirmada pelos documentos neotestamentários. ``Cristo ressuscitou dos mortos! Morrendo, destruiu a morte; aos mortos, deu a vida.'' É evento histórico e transcendente, pregado como parte essencial do Mistério Pascal junto à Cruz.

O túmulo vazio, descoberto pelas santas mulheres e apóstolos, é sinal inicial: ``Por que buscais entre os mortos aquele que vive? Ele não está aqui, ressuscitou''. Não prova direta, mas essencial, refutando roubo ou retorno à vida terrena como Lázaro.

As aparições do Ressuscitado confirmam: a Pedro, aos Doze, aos discípulos de Emaús, a Tomé, aos 500, a Paulo. Ele come, é tocado, mas entra glorificado, além do espaço-tempo, enchendo-se do Espírito Santo.

A Igreja confessa: Cristo não reviveu, mas passou à vida divina, primeiro fruto da ressurreição geral. Contra gnósticos e modernistas, é fato público, fundamento da pregação apostólica.

Para adultos a partir dos 20 anos no RICA, meditar a Ressurreição no tempo pascal da mistagogia fortalece a fé batismal recém-recebida, inspirando testemunho vivo do Ressuscitado na comunidade neófita.

\section{A Ascensão de Jesus ao Céu (\S 659--667)}

A Ascensão marca a entrada definitiva da humanidade de Jesus no domínio celestial de Deus, donde voltará. Após 40 dias de aparições, sobe aos olhos dos Apóstolos no Monte das Oliveiras, nuvem o ocultando como no Tabor.

Não ausência, mas exaltação gloriosa: corpo glorioso entra no céu, centro da liturgia celeste. ``Elevado da terra, atrairá todos a si'', unindo cruz e céu em movimento pascal contínuo.

Cristo, sumo sacerdote eterno, intercede por nós ante o Pai, principal ator da liturgia eucarística. Reina como Rei-sacerdote, submetendo tudo ao Pai.

Esse mistério realiza a Encarnação: humanidade assumida eleva-se à glória trinitária, prometendo nossa participação.

Na catequese para adultos não iniciados, a Ascensão no rito de envio pós-batismo eleva aspirações celestiais, motivando oração litúrgica e missão evangelizadora durante o catecumenato.

\section{Está Sentado à Direita de Deus Pai Todo-Poderoso (\S 668--682)}

Porém, reino é já e não ainda: luta contra o Anticristo precede consumação. Cristo submete inimigos, inclusive a morte, até o fim dos tempos.

Seu senhorio é universal e escatológico: julgará vivos e mortos, separando justos como trigo do joelho. A Igreja peregrina participa de sua realeza batismal.

Contra milenarismo, não reino terreno milenar; sim, vitória final na Parusia, renovando cosmos.

Para catecúmenos adultos, professar o reinado de Cristo no rito de eleição dissipa utopias seculares, ancorando esperança na liturgia dominical e obras de justiça social rumo à Crisma régia.

\section{Conclusão}

Professar Ressurreição, Ascensão e reinado de Cristo proclama vitória sobre morte e mal. Adultos no RICA, vivam essa glória na iniciação pascal, reinando com Ele.